This paper introduced HP-PDHG, a hybrid classical primal heuristic for converting strong relaxation trajectories into feasible mixed-integer incumbents. The method combines population PDHG, a barrier-aware non-local perturbation, progressive integrality projection, conservative repair, and an optional two-phase controller. Its contribution is not a claim about quantum hardware, but an interpretable search architecture for feasibility recovery in MIP.

The audited experiments support this positioning at three levels. First, on the four-instance feasibility-recovery benchmark, the baseline produces attractive but infeasible candidates on every case, whereas HP-PDHG recovers strict feasibility on all four; across 80 paired runs, the proposed solver is feasible in 80/80 versus 0/80 for the baseline. Second, on a 12-instance official MIPLIB 2017 extension, HP-PDHG lowers mean audited violation on eight instances, ties on two, loses on two, and wins a paired Wilcoxon test on per-instance mean violations. Third, on a representative six-instance population-based comparison against GA, DE, SHADE, and SL-PSO, the method attains the best average Friedman rank. Ablation shows that progressive integrality projection and repair drive the gains, while the barrier-aware perturbation mainly broadens exploration. The 390-run sensitivity analysis further indicates a broad practical plateau rather than a brittle hyperparameter knife-edge.

\subsection{Interpretation}
\label{subsec:future}

The main lesson is that feasibility-oriented heuristic design deserves to be treated as a first-class scientific problem rather than as post-processing attached to LP optimization. For MIP, the difference between an excellent but infeasible point and a valid incumbent is operationally decisive. HP-PDHG is valuable because it treats that transition explicitly: PDHG handles the relaxation, integrality projection contracts the discrete residual, repair enforces admissibility, and the barrier-aware perturbation helps the population escape repeated infeasible basins.

The larger-scale extensions make this interpretation more precise. The official MIPLIB 2017 study shows that the benefit survives beyond the original four audited cases, but also makes clear that the present implementation is still stronger at \emph{reducing} infeasibility than at \emph{closing} it on every harder instance under a modest budget. The population-based baseline study reaches a similar conclusion: the method is well placed among standard intelligent heuristics, but the evidence does not justify pairwise universal-superiority claims.

\subsection{Limitations and Future Directions}
\label{subsec:impact}

Several limitations remain. Even after the official extension, the strict-feasibility benchmark is still modest in size, the method is stronger at \emph{reducing} distance to feasibility than at \emph{consistently closing} it on every medium-scale instance, some representative cases still favor conventional metaheuristics, and the current implementation remains predominantly CPU based despite obvious opportunities for acceleration.

These limitations define several next steps. The most immediate is to embed HP-PDHG as an incumbent-generation component inside a modern exact solver and measure its downstream effect on pruning and total solve time. Other priorities are to enlarge the official benchmark study, add learned guidance for perturbation or projection scheduling, and explore adaptive portfolio control that decides when relaxation-centered search should dominate and when a more generic intelligent operator should take over. The same design may also extend beyond linear MIP if the energy function and repair model are adapted carefully.

More broadly, the paper argues for a disciplined interaction between ideas from quantum optimization, operations research, and intelligent heuristic design. The most productive transfer may come not from claims about speedup, but from extracting algorithmic principles that withstand classical scrutiny: barrier-sensitive exploration, staged commitment to discrete structure, and explicit feasibility preservation.
